\section*{Capítulo \ref{ch:b3:milky-way-galaxies} --- La Vía Láctea y las galaxias}

\begin{solution}{exo:b3:milky-way-galaxies:1}
(a)~$2\pi R_0 = 2\pi\times\qty{2.523e17}{km} = \qty{1.585e18}{km}$; a
\qty{229}{km.s^{-1}} esto requiere \qty{6.92e15}{s}, es decir
\qty{219}{Myr}. (b)~$4540/219 = 20.7$: unas veinte vueltas.
\end{solution}

\begin{solution}{exo:b3:milky-way-galaxies:2}
(a)~$\qty{1}{au}/\qty{1}{yr} = \num{1.496e8}/\num{3.156e7} =
\qty{4.74}{km.s^{-1}}$. (b)~$v = 4.74\times6.411\times8.178 =
\qty{248.5}{km.s^{-1}}$, dentro del valor de la revisión
\qty[separate-uncertainty=true]{248 \pm 3}{km.s^{-1}}. (c)~$248.5 - 11 = \qty{237.5}{km.s^{-1}}$,
unos \qty{8}{km.s^{-1}} por encima de los \qty{229}{km.s^{-1}} de las gigantes rojas:
los dos se apoyan en valores distintos de $R_0$ (8.178 y \qty{8.122}{kpc}),
y el propio \omterm{def:b3:milky-way-galaxies:lsr}{estándar local de reposo} puede desplazarse hasta unos \qty{15}{km.s^{-1}} respecto
a una órbita verdaderamente circular.
\end{solution}

\begin{solution}{exo:b3:milky-way-galaxies:3}
(a)~$n = 10(1 - 0.6) = 4$: E4. (b)~$b/a = 1 - 7/10 = 0.3$. (c)~La clasificación
usa la razón aparente de ejes: un elipsoide achatado visto a lo largo de su eje corto
tiene un contorno circular y se clasifica como E0.
\end{solution}

\begin{solution}{exo:b3:milky-way-galaxies:4}
$d = 10^{\mu/5 + 1}$~pc: la LMC \qty{49.6}{kpc}, Andrómeda
\qty{783}{kpc}, NGC~4258 \qty{7.58}{Mpc}. Una estrella de la Nube tiene una
paralaje de $1/\qty{49.6}{kpc} = \qty{0.020}{mas}$, del orden de la incertidumbre del mejor
catálogo para estrellas brillantes: una sola paralaje no da ninguna
distancia, y la Nube se sitúa por otros medios (binarias de eclipse),
a partir de las cuales arranca la escala de las \omterm{def:b3:milky-way-galaxies:cepheid}{variables cefeidas}.
\end{solution}

\begin{solution}{exo:b3:milky-way-galaxies:5}
(a)~$\omega_0 = A - B = \qty{27.2}{km.s^{-1}.kpc^{-1}}$ y $V_0 =
27.2\times8.178 = \qty{222}{km.s^{-1}}$. (b)~$V'_0 = -(A + B) =
\qty{-3.4}{km.s^{-1}.kpc^{-1}}$: una curva que desciende suavemente. (c)~Ambas están
por debajo de los valores de las gigantes rojas (229 y \qty{-1.7}{km.s^{-1}} por
kiloparsec); los autores de esa curva hallaron lo mismo: un $V_0$ unas dos
desviaciones típicas más bajo a causa de $R_0$ y de estas \omterm{def:b3:milky-way-galaxies:rotation-curve}{constantes de Oort}. Las
fórmulas suponen órbitas circulares en un disco axisimétrico; los valores no nulos de $C$ y
$K$ muestran que el campo de velocidades local está perturbado por la barra y los
brazos espirales, y las $A$ y $B$ locales incorporan esas perturbaciones.
\end{solution}

\begin{solution}{exo:b3:milky-way-galaxies:6}
(a)~$\rho = V^2/(4\pi GR^2) = 229^2/(4\pi\times\num{4.301e-6}\times
8.178^2) = \num{1.45e7}$ masas solares por kiloparsec cúbico, o
\num{0.0145} por parsec cúbico. (b)~$M_\odot c^2 = \num{1.9884e30}\times
(\num{2.998e8})^2 = \qty{1.787e47}{J} = \qty{1.115e57}{GeV}$, y
$\qty{1}{pc^3} = (\qty{3.0857e18}{cm})^3 = \qty{2.938e55}{cm^3}$: la
razón es \qty{37.96}{GeV.cm^{-3}}. Entonces $0.30/37.96 = \num{0.0079}$ masas solares
por parsec cúbico. (c)~La densidad asociada a la curva plana cuenta toda la masa
que sostiene las órbitas, incluidas las estrellas y el gas, como si estuviera repartida en una
esfera; la materia oscura es solo una parte de ella, aquí alrededor de la mitad.
\end{solution}

\begin{solution}{exo:b3:milky-way-galaxies:7}
(a)~$W_I(\mathrm{LMC}) = -3.314\times1.4771 + 15.888 = \qty{10.993}{mag}$,
así $M_W = 10.993 - 18.477 = \qty{-7.484}{mag}$. (b)~$\mu = 25.00 + 7.484 =
32.484$ y $d = 10^{32.484/5 + 1}~\mathrm{pc} = \qty{31.4}{Mpc}$.
(c)~$\delta d/d = (\ln 10/5)\,\delta\mu = 0.4605\times0.077 =
\qty{3.5}{\percent}$ para una estrella, $3.5/\sqrt{50} = \qty{0.5}{\percent}$
para cincuenta, si sus errores son independientes (el error de la distancia propia de la Nube
se suma por igual a todas ellas).
\end{solution}

\begin{solution}{exo:b3:milky-way-galaxies:8}
$\mu = I - A_I - M_I = 26.00 - 0.05 + 4.05 = 30.00$, así $d =
10^{7}~\mathrm{pc} = \qty{10.0}{Mpc}$. Un error de \qty{0.04}{mag} equivale a
$0.4605\times0.04 = \qty{1.8}{\percent}$ en distancia.
\end{solution}

\begin{solution}{exo:b3:milky-way-galaxies:9}
(a)~$d = cz/H_0$ da $5\log_{10}d = 5\log_{10}(cz) - 5\log_{10}H_0$,
así $m_B - M_B - 25 = 5\log_{10}(cz) - 5\log_{10}H_0$; con
$\log_{10}(cz) = a_B + 0.2\,m_B$, $5\log_{10}H_0 = 5a_B + M_B + 25$, lo que
es el resultado. (b)~$\log_{10}H_0 = -3.8506 + 0.714158 + 5 = 1.8636$, así
$H_0 = \qty{73.04}{km.s^{-1}.Mpc^{-1}}$. (c)~$M_B = 5(\log_{10}67.4 -
5.714158) = \qty{-19.427}{mag}$: supernovas más brillantes en \qty{0.174}{mag},
por tanto, todas las distancias son mayores en un factor $10^{0.2\times0.174} = 1.084$, la
razón $73.04/67.4$. La tensión es una discrepancia de un ocho por ciento en la
escala de distancias.
\end{solution}

\begin{solution}{exo:b3:milky-way-galaxies:10}
(a)~Una estrella en el borde, en una órbita circular, tiene $V^2/R \approx GM/R^2$
(aproximación esférica). (b)~$R = \sqrt{M/(\pi\Sigma)}$, así $V^2 =
GM/R = G\sqrt{\pi\Sigma}\,M^{1/2}$ y $M = V^4/(\pi\Sigma G^2) \propto
V^4$; con una razón masa-luminosidad común, $L \propto V^4$ también.
(c)~$(100/229)^4 = \num{0.036}$: la galaxia más rápida es 27 veces más
masiva. (d)~Con igual luminosidad, el flujo cae como $1/d^2$, así que la más
débil está al doble de distancia.
\end{solution}

\begin{solution}{exo:b3:milky-way-galaxies:11}
(a)~$\qty{1}{km.s^{-1}} = \qty{1.0227}{kpc.Gyr^{-1}}$, así
\qty{229}{km.s^{-1}} equivale a \qty{234.2}{kpc.Gyr^{-1}}. (b)~$2\pi R/V$:
\qty{107}{Myr} a \qty{4}{kpc}, \qty{322}{Myr} a \qty{12}{kpc}.
(c)~$\Delta\phi = Vt\,(1/4 - 1/12) = 234.2\times0.1667 = \qty{39.0}{rad}$,
6.2 vueltas: el brazo sería una espiral de seis vueltas apretadas entre 4 y
\qty{12}{kpc}, mucho más cerrada que los brazos observados, y aún más cerrada en cada
nueva rotación. (d)~En una onda de densidad el patrón gira a una única velocidad angular
como un sólido rígido; las estrellas y el gas lo atraviesan, así que la forma
de los brazos no se enrolla.
\end{solution}

\begin{solution}{exo:b3:milky-way-galaxies:12}
(a)~$a = (4.2\times10^6\times16^2)^{1/3} = \qty{1024}{au}$; a
\qty{8178}{pc} subtiende $1024/8178 = \ang{;;0.125}$ (una ua a un parsec
subtiende un arcosegundo). (b)~$R_S = 2GM/c^2 = 2\times\num{6.674e-11}
\times\num{8.35e36}/(\num{2.998e8})^2 = \qty{1.24e10}{m} = \qty{0.083}{au}$;
el diámetro $2R_S$ subtiende $0.166/8178 = \num{2.0e-5}$ arcosegundos, unos
\qty{20}{\micro as}. (c)~$R_S = \qty{4.02e-7}{pc}$, y $M/(\tfrac43\pi
R_S^3) = \num{1.5e25}$ masas solares por parsec cúbico, cuarenta veces el
límite inferior fijado por el movimiento y el tamaño de Sgr~A*: el límite queda a
menos de un factor cuarenta de la densidad del propio agujero negro.
\end{solution}

\begin{solution}{pb:b3:milky-way-galaxies:1}

\textbf{Parte I.}

\textbf{1.} $v_r = A\,d\sin2\ell$: $+A d$ en \ang{45}, 0 en \ang{90}, $-A d$ en
\ang{135}, 0 en \ang{180}.

\textbf{2.} $A = (15.4 + 15.2)/(2\times\qty{1.00}{kpc}) =
\qty{15.3}{km.s^{-1}.kpc^{-1}}$.

\textbf{3.} En \ang{90}, $v_t = d(B - A) = 4.74\times(-5.91)\times1.00 =
\qty{-28.02}{km.s^{-1}}$, así que $B = 15.3 - 28.02 =
\qty{-12.72}{km.s^{-1}.kpc^{-1}}$.

\textbf{4.} En \ang{0}, $v_t = d(A + B) = 4.74\times0.55 =
\qty{2.61}{km.s^{-1}}$, así que $B = 2.61 - 15.3 = \num{-12.69}$: los dos
coinciden dentro de \num{0.03}.

\textbf{5.} $\omega_0 = A - B = \qty{28.0}{km.s^{-1}.kpc^{-1}}$ y
$V_0 = 28.0\times8.18 = \qty{229}{km.s^{-1}}$.

\textbf{6.} $V'_0 = -(A + B) = \qty{-2.6}{km.s^{-1}.kpc^{-1}}$: disminuye suavemente.

\textbf{7.} $P = 2\pi/\omega_0$: $2\pi\times\qty{8.18}{kpc}/\qty{229.2}{km.s^{-1}}
= \qty{219}{Myr}$; $4540/219 = 20.7$ vueltas.

\textbf{Parte II.}

\textbf{8.} El punto tangente requiere una línea de visión que pase más cerca del
centro que el Sol, $\sin\ell > 0$ con $\ell < \ang{90}$ (o el
cuadrante simétrico, $\ang{270} < \ell < \ang{360}$, donde la mayor
velocidad es la más negativa). Fuera del círculo solar la distancia
al centro crece a lo largo de toda la línea de visión: ningún punto destaca,
y la velocidad por sí sola no indica dónde se encuentra el gas.

\textbf{9.} $R = 8.18\sin\ell$: \qtylist{5.26;6.27;7.08;7.69}{kpc}.

\textbf{10.} $V = v_\mathrm{max} + V_0\sin\ell$: $82.9 + 147.2 = 230.1$,
$55.6 + 175.4 = 231.0$, $32.4 + 198.3 = 230.7$, $14.9 + 215.2 =
\qty{230.1}{km.s^{-1}}$.

\textbf{11.} $\Delta\nu = \nu v/c = \num{1.420406e9}\times82.9/299792 =
\qty{393}{kHz}$, hacia frecuencias más bajas: el gas se aleja.

\textbf{12.} Dentro de unos \qty{3}{km.s^{-1}} de las gigantes rojas: entre 5 y
\qty{8}{kpc} la curva es plana.

\textbf{13.} Una masa puntual da $V = V_0\sqrt{R_0/R} = 229\times\sqrt{8.18/5.26} =
\qty{286}{km.s^{-1}}$ en el primer punto tangente, muy superior a la
\qty{230}{km.s^{-1}} hallada: la masa no está concentrada en el centro,
buena parte de ella se encuentra entre 5 y \qty{8}{kpc} y más allá.

\textbf{Parte III.}

\textbf{14.} $\num{1.067e10} + 2\times\num{3.944e10} = \num{8.96e10}$ masas
solares, casi el doble de la cifra de la revisión \num[separate-uncertainty=true]{5 \pm 1 e10}; las estrellas del modelo
pueden ser demasiado pesadas, lo que haría que la materia oscura hallada más abajo
sea una subestimación.

\textbf{15.} $V_b^2 = GMR^2/(R^2 + b^2)^{3/2} = \num{4.301e-6}\times\num{1.067e10}
\times400/(400.09)^{3/2}$, $V_b = \qty{47.9}{km.s^{-1}}$.

\textbf{16.} Disco delgado: $(a + b)^2 = 30.80$, $V^2 = \num{4.301e-6}\times\num{3.944e10}
\times400/(430.8)^{3/2}$, $V = \qty{87.1}{km.s^{-1}}$. \omterm{def:b3:milky-way-galaxies:structure}{Disco grueso}:
$(a + b)^2 = 11.56$, $V = \qty{90.1}{km.s^{-1}}$.

\textbf{17.} $V_\mathrm{vis} = \sqrt{47.9^2 + 87.1^2 + 90.1^2} =
\qty{134.2}{km.s^{-1}}$.

\textbf{18.} A \qty{8.18}{kpc}: 74.8, 108.4 y \qty{127.8}{km.s^{-1}},
en conjunto \qty{183.5}{km.s^{-1}}, frente a $V_0 = \qty{229}{km.s^{-1}}$:
la masa visible resulta insuficiente ya en el Sol.

\textbf{19.} $\sqrt{G\times\num{8.96e10}/20} = \qty{138.8}{km.s^{-1}}$. La
velocidad del modelo, \qty{134.2}{km.s^{-1}}, es cercana, porque a \qty{20}{kpc},
entre cuatro y seis longitudes de escala del disco hacia fuera, casi toda la masa visible se encuentra
dentro de la órbita: la curva visible se ha vuelto kepleriana.

\textbf{Parte IV.}

\textbf{20.} $M = 20\times200^2/\num{4.301e-6} = \num{1.860e11}$ masas solares.

\textbf{21.} $M_\mathrm{vis} = 20\times134.2^2/\num{4.301e-6} = \num{8.37e10}$
masas solares.

\textbf{22.} $M_\mathrm{dark} = \num{1.860e11} - \num{8.37e10} = \num{1.02e11}$;
ya que ambas masas son $RV^2/G$ al mismo $R$, la fracción es
$1 - (V_\mathrm{vis}/V)^2 = 1 - (134.2/200)^2 = \num{0.550}$.

\textbf{23.} $x = 20/14.8 = 1.351$, $\ln(2.351) - 1.351/2.351 = 0.8549 - 0.5748
= 0.2801$, $M = 4\pi\times\num{1.06e7}\times14.8^3\times0.2801 =
\num{1.21e11}$ masas solares; la fracción oscura es
$1.21/(1.21 + 0.837) = \num{0.59}$. El halo se ajusta a toda la
curva, que a \qty{20}{kpc} la sobrepasa: las medidas caen allí
unos \qty{10}{km.s^{-1}} por debajo del modelo suave (\qty{209}{km.s^{-1}}).

\textbf{24.} $x = 8.18/14.8 = 0.553$, $\rho = \num{1.06e7}/(0.553\times1.553^2)
= \num{7.95e6}$ masas solares por kiloparsec cúbico, \num{0.0080} por
parsec cúbico, o \qty{0.30}{GeV.cm^{-3}}, el valor que citan los autores del
ajuste.

\textbf{25.} $1 - 8.37/18.60 = \num{0.550}$: la fracción de materia oscura de la
masa dentro de \qty{20}{kpc} es
\textbf{\qty[text-series-to-math=true]{55}{\percent}}.
\end{solution}
